% Note: Several packages (booktabs, xcolor, colortbl, graphicx, caption,
% enumitem, mdframed, microtype, amssymb, amsmath, amsfonts, pifont,
% tabularx, tcolorbox, authblk, hyperref, url) are already loaded by shengshu.cls.
% Do not reload them here.

\usepackage[utf8]{inputenc}
\usepackage[T1]{fontenc}
\usepackage{charter}

% Match lab.cls bold font sizing
\newcommand{\bffont}{\fontsize{8.8}{10}\selectfont}
\DeclareTextFontCommand{\textbf}{\bffont\bfseries\selectfont}

\usepackage{wrapfig}
\usepackage{xspace}
\usepackage{textgreek}
\usepackage{bm}
\usepackage{cleveref}
\usepackage{multirow}
\usepackage{subcaption}
\usepackage{orcidlink}
\usepackage{footmisc}
\usepackage{algorithm}
\usepackage{algpseudocode}
\usepackage{listings}
\usepackage{csquotes}
\usepackage{marvosym}
\usepackage{nicefrac}
\usepackage{multicol}
\usepackage{tabto}
\usepackage{adjustbox}
\usepackage{dblfloatfix}
\usepackage{verbatim}
\usepackage{mathtools}
\usepackage{array}
\usepackage{bbm}
\usepackage{makecell}
\usepackage{siunitx}
\usepackage{pdflscape}
\usepackage{arydshln}
\usepackage{hhline}
\usepackage{diagbox}
\usepackage[square,sort,comma,numbers]{natbib}
\usepackage{fp}

\def\UrlBreaks{\do\/\do-}

% Project colors
\definecolor{linkc}{rgb}{0, 0.44, 0.74}
\definecolor{eqc}{rgb}{1, 0, 0}
\definecolor{newcitecolor}{rgb}{0,0.6,0}
\definecolor{mygreen}{RGB}{34,139,34}
\definecolor{mylightblue}{RGB}{0,162,230}
\definecolor{deepyellow}{RGB}{255, 215, 0}
\definecolor{catgray}{gray}{0.92}
\definecolor{pearDark}{RGB}{171, 195, 87}
\definecolor{codebg}{RGB}{245, 245, 245}
\definecolor{keywordcolor}{RGB}{0, 0, 153}
\definecolor{commentcolor}{RGB}{34, 139, 34}
\definecolor{stringcolor}{RGB}{163, 21, 21}
\definecolor{numbercolor}{RGB}{128, 128, 128}

\newcommand{\xmarker}{\ding{55}}
\newcommand{\cmarker}{\ding{51}}
\newcommand{\sota}[1]{\cellcolor{pearDark!20}{#1}}
\newcommand{\change}[1]{\textcolor{orange}{#1}}

% Co-author comment macros
\newcommand{\zm}[1]{{\color{red}{[[\textbf{zm: }#1]]}}}
\newcommand{\hz}[1]{{\color{red}{[[\textbf{hz: }#1]]}}}

% Method name shortcut
\newcommand{\ours}{Causal Forcing++\xspace}

% Common math shortcuts
\newcommand{\bx}{\boldsymbol{x}}
\newcommand{\bz}{\boldsymbol{z}}

% \onedot (from CVPR style) for abbreviation periods
\def\onedot{\futurelet\@let@token\@onedot}
\def\@onedot{\ifx\@let@token.\else.\null\fi\xspace}

\def\eg{\emph{e.g}\onedot} \def\Eg{\emph{E.g}\onedot}
\def\ie{\emph{i.e}\onedot} \def\Ie{\emph{I.e}\onedot}
\def\cf{\emph{cf}\onedot} \def\Cf{\emph{Cf}\onedot}
\def\etc{\emph{etc}\onedot} \def\vs{\emph{vs}\onedot}
\def\wrt{w.r.t\onedot} \def\dof{d.o.f\onedot}
\def\iid{i.i.d\onedot} \def\wolog{w.l.o.g\onedot}
\def\etal{\emph{et al}\onedot}
\makeatletter
\def\blfootnote#1{\xdef\@thefnmark{}\@footnotetext{\scriptsize #1}}
\makeatother

\hypersetup{
    colorlinks=true,
    urlcolor=shengshublue,
}

\let\cite\citep

\lstset{
    language=Python,
    basicstyle=\ttfamily\footnotesize,
    backgroundcolor=\color{codebg},
    keywordstyle=\color{keywordcolor}\bfseries,
    commentstyle=\color{commentcolor}\itshape,
    stringstyle=\color{stringcolor},
    numbers=none,
    numberstyle=\color{numbercolor}\tiny,
    stepnumber=1,
    numbersep=5pt,
    showspaces=false,
    showstringspaces=false,
    breaklines=true,
    frame=none,
    framesep=8pt,
    framerule=0.5pt
}
